\documentclass[12pt,a4paper]{article}

% --- Packages ---
\usepackage[utf8]{inputenc}
\usepackage[T1]{fontenc}
\usepackage[indonesian]{babel}
\usepackage{geometry}
\geometry{top=2.5cm, bottom=2.5cm, left=2.5cm, right=2.5cm}
\usepackage{hyperref}
\usepackage{listings}
\usepackage{xcolor}
\usepackage{booktabs}
\usepackage{titlesec}
\usepackage{enumitem}

% --- Syntax Highlighting for TypeScript ---
\definecolor{codegreen}{rgb}{0,0.6,0}
\definecolor{codegray}{rgb}{0.5,0.5,0.5}
\definecolor{codepurple}{rgb}{0.58,0,0.82}
\definecolor{backcolour}{rgb}{0.95,0.95,0.92}

\lstdefinelanguage{TypeScript}{
  keywords={abstract, any, as, boolean, break, case, catch, class, console, const, continue, debugger, declare, default, delete, do, else, enum, export, extends, false, finally, for, from, function, get, if, implements, import, in, instanceof, interface, let, module, never, new, null, number, object, package, private, protected, public, readonly, require, return, set, static, string, super, switch, this, throw, true, try, type, typeof, var, void, while, with, yield, async, await},
  sensitive=true,
  comment=[l]{//},
  morecomment=[s]{/*}{*/},
  morestring=[b]',
  morestring=[b]"
}

\lstset{
    language=TypeScript,
    backgroundcolor=\color{backcolour},   
    commentstyle=\color{codegreen},
    keywordstyle=\color{magenta},
    numberstyle=\tiny\color{codegray},
    stringstyle=\color{codepurple},
    basicstyle=\ttfamily\small,
    breakatwhitespace=false,         
    breaklines=true,                 
    captionpos=b,                    
    keepspaces=true,                 
    numbers=left,                    
    numbersep=5pt,                  
    showspaces=false,                
    showstringspaces=false,
    showtabs=false,                  
    tabsize=2
}

% --- Title ---
\title{
    \textbf{Laporan Pengembangan Aplikasi}\\
    \large \textit{Mindful Pulse: Student Wellness Tracker}
}
\author{Dokumentasi Teknis Otomatis}
\date{\today}

\begin{document}

\maketitle
\tableofcontents
\newpage

\section{Deskripsi Aplikasi}
\textbf{Mindful Pulse} adalah aplikasi \textit{wellness-tracker} berbasis web yang dirancang khusus untuk mahasiswa guna mengelola kesehatan fisik dan mental di tengah kesibukan akademik.

\begin{itemize}[label=--]
    \item \textbf{Tujuan}: Membantu mahasiswa menjaga keseimbangan antara performa akademik dan kesejahteraan diri melalui pelacakan kebiasaan dan dukungan AI.
    \item \textbf{Target Pengguna}: Mahasiswa universitas yang menghadapi stres akademik, manajemen waktu yang buruk, atau pengabaian kesehatan fisik.
    \item \textbf{Masalah yang Diselesaikan}: Kurangnya kesadaran akan hidrasi, pola tidur, manajemen tugas yang membebani, serta kebutuhan akan motivasi yang personal.
    \item \textbf{Fitur Unggulan}:
    \begin{itemize}
        \item \textbf{The Daily 5}: Dashboard ringkas untuk memantau 5 memtrik kesehatan harian.
        \item \textbf{AI Wellness Coach}: Asisten bertenaga Gemini yang memberikan rekomendasi berdasarkan data perilaku pengguna.
        \item \textbf{Gamification}: Sistem \textit{streak} dan \textit{achievements} untuk memotivasi konsistensi pengguna.
    \end{itemize}
\end{itemize}

\section{Fungsi-Fungsi Utama}

\subsection{Fungsi Water Tracker}
\begin{itemize}
    \item \textbf{Tambah Air}: Pengguna dapat menambah jumlah asupan air (ml atau L).
    \item \textbf{Progress \%}: Menampilkan persentase pencapaian terhadap target harian (default 3L).
    \item \textbf{Mode Puasa}: Penyesuaian pengingat dan target saat periode puasa.
\end{itemize}

\subsection{Fungsi Mood Tracker}
\begin{itemize}
    \item \textbf{Pilih Mood}: Input skala perasaan 1-4.
    \item \textbf{Simpan Jurnal}: Catatan teks untuk refleksi diri.
    \item \textbf{Riwayat}: Melihat tren perasaan selama seminggu terakhir untuk deteksi dini stres.
\end{itemize}

\subsection{Fungsi Jadwal Kuliah}
\begin{itemize}
    \item \textbf{Tambah Jadwal}: Menyimpan data mata kuliah, ruang, dan jam.
    \item \textbf{Tampil Hari Ini}: Secara dinamis menampilkan jadwal yang relevan dengan hari berjalan di dashboard.
\end{itemize}

\subsection{Fungsi AI Insight \& Analytics}
\begin{itemize}
    \item \textbf{Rekomendasi Otomatis}: AI menganalisis data perilaku dan memberikan motivasi instan.
    \item \textbf{Breakdown Tugas}: Menguraikan tugas kuliah yang besar menjadi sub-tugas (micro-tasks).
    \item \textbf{Grafik Mingguan}: Visualisasi tren data kesehatan selama 7 hari.
\end{itemize}

\section{Dokumentasi OOP \& FP}

\subsection{Object-Oriented Programming (OOP)}
Aplikasi menggunakan pola OOP untuk layanan (services) guna mengenkapsulasi logika bisnis yang kompleks dan menjaga \textit{single source of truth}.

\begin{lstlisting}[caption=Implementasi Singleton Pattern pada AIService]
class AIService {
  private static instance: AIService;
  private lastAnalysisTime: number = 0;

  public static getInstance(): AIService {
    if (!AIService.instance) {
      AIService.instance = new AIService();
    }
    return AIService.instance;
  }

  public async analyzeWellness(force: boolean = false): Promise<AIAnalysisResult | null> {
    const context = this.generateContext();
    try {
      const response = await fetch("/api/ai/analyze-wellness", {
        method: "POST",
        body: JSON.stringify({ context }),
      });
      return await response.json();
    } catch (error) {
      this.handleError(error, "Wellness Analysis");
      return null;
    }
  }
}
\end{lstlisting}

\subsection{Functional Programming (FP)}
Logika kalkulasi di dalam store dan utility menggunakan prinsip FP agar bersifat \textit{pure} dan \textit{immutable}.

\begin{lstlisting}[caption=Contoh Pure Function untuk Kalkulasi Skor]
export const calculateDailyWellnessScore = (state: WellnessState, date: string): number => {
  const water = state.waterLogs[date] || 0;
  const steps = state.stepsLogs[date] || 0;

  const weights = { water: 0.5, steps: 0.5 };
  const scores = {
    water: Math.min(water / state.baseWaterGoal, 1),
    steps: Math.min(steps / state.stepGoal, 1),
  };

  return Math.round((scores.water * weights.water + scores.steps * weights.steps) * 100);
};
\end{lstlisting}

\section{Error Handling}
Aplikasi menerapkan sistem pertahanan berlapis:
\begin{enumerate}
    \item \textbf{Validasi Store}: Menggunakan \texttt{Math.max} untuk mencegah nilai negatif.
    \item \textbf{Data Kosong}: Penggunaan \textit{optional chaining} (\texttt{?.}) dan nilai default.
    \item \textbf{API Throttling}: Melindungi kuota API Gemini dari pemanggilan berlebihan.
\end{enumerate}

\begin{lstlisting}[caption=Implementasi Try-Catch dengan Fallback Data]
try {
  const response = await fetch("/api/tasks/breakdown", {
    method: "POST",
    body: JSON.stringify({ taskTitle }),
  });
  if (!response.ok) throw new Error("Service unavailable.");
  const data = await response.json();
  return data.subtasks;
} catch (error) {
  console.error("Critical Failure", error);
  return ["Persiapan", "Pengerjaan", "Finalisasi"]; // Fallback
}
\end{lstlisting}

\section{Teknik Pengujian}
\begin{itemize}
    \item \textbf{Unit Test}: Menguji fungsi kalkulasi skor secara terisolasi.
    \item \textbf{Integration Test}: Sinkronisasi antar fitur (misal: Input Air ke Dashboard).
    \item \textbf{UI Test}: Responsivitas pada \textit{viewport} perangkat mobile.
\end{itemize}

\begin{table}[h!]
\centering
\caption{Hasil Pengujian Sistem}
\vspace{2mm}
\begin{tabular}{@{}llll@{}}
\toprule
\textbf{Fitur} & \textbf{Jenis Test} & \textbf{Kondisi} & \textbf{Status} \\ \midrule
Water Tracker & Unit & Input -250ml & PASS \\
Mood Tracker & Integration & Simpan Mood $\rightarrow$ Update Streak & PASS \\
AI Insight & Integration & API Timeout & PASS \\
Dashboard & UI & iPhone SE Viewport & PASS \\
Task List & Edge Case & Judul $>$ 500 Karakter & PASS \\ \bottomrule
\end{tabular}
\end{table}

\vfill
\begin{center}
\textit{Laporan ini dihasilkan secara otomatis sebagai dokumentasi teknis Mindful Pulse.}
\end{center}

\end{document}
